\chapter{Vertebroplasty}

\section*{Overview}
Vertebroplasty is a minimally invasive procedure where medical-grade bone cement (PMMA) is injected into a fractured vertebra to stabilize it and relieve pain. The exact approach, setting, and preparation depend on the indication, anatomy, and the patient's health.

\section*{Alternative Names}
Spinal cement injection, Percutaneous vertebroplasty

\section*{Principle}
Bone cement is injected under imaging guidance through a needle directly into a fractured vertebral body. The cement hardens within minutes, stabilizing the vertebra and reducing the pain caused by movement of the fracture fragments.
\section*{History}
First performed in France in 1984 by Galibert and Deramond to treat an aggressive vertebral hemangioma.

\section*{Why It Is Done}
Ordered for painful vertebral compression fractures caused by osteoporosis or bone metastasis that do not respond to conservative care.

\section*{Is This a Primary or Secondary Test or Procedure}
Secondary treatment option for painful vertebral compression fractures when conservative treatment fails.

\section*{How It Is Performed}
Under local anesthesia and conscious sedation. Under continuous fluoroscopic guidance, a needle is placed into the vertebral body, and liquid PMMA is injected.

\section*{Preparation}
The team reviews bleeding risk and medicines, including anticoagulants and antiplatelet drugs; do not stop them without a clinician-directed plan. MRI or another study may be used when clinically indicated to assess fracture acuity and exclude other causes of pain.

\section*{Contraindications and Precautions}
Active systemic or local osteomyelitis; uncorrected coagulopathy; fracture causing severe spinal cord compression.

\section*{Risks}
Cement leakage (into spinal canal or veins), pulmonary cement embolism, nerve root compression, or adjacent level fractures.

\section*{Aftercare and Recovery}
Lie flat for 1-2 hours until cement hardens; monitor vital signs; assess pain relief (often immediate).

\section*{Outcome and Follow-up}
Success is marked by immediate reduction in back pain and stabilized vertebral height.

\section*{Expected Results}
Stable vertebral body with uniform cement distribution; significant reduction in pain; no cement leakage.

\section*{Follow-up}
Re-evaluation in 1-2 weeks; bone density management (osteoporosis therapy).

\section*{When to Seek Medical Care}
Follow the procedure team's specific instructions. Seek urgent medical assessment for severe or worsening pain, uncontrolled bleeding, fever or signs of infection, trouble breathing, chest pain, fainting, new weakness or confusion, inability to urinate, or any rapidly worsening symptom. The appropriate warning signs depend on the procedure and the patient's medical condition.

\section*{References}
\begin{enumerate}
  \item MedlinePlus. Vertebroplasty. U.S. National Library of Medicine; 2025.
  \item Current Medical Diagnosis and Treatment. McGraw-Hill; 2025.
  \item American Society of Regional Anesthesia and Pain Medicine. Antiplatelet and anticoagulant medication guidance for interventional spine procedures. \url{https://asra.com/the-asra-family/special-interest-groups/neuromodulation-sig/neuromodulation-sig-article-item/guidelines/2022/12/14/interventional-spine-and-pain-procedures-in-patients-on-antiplatelet-and-anticoagulant-medications-%28second-edition%29} (accessed 2026-08-16).
\end{enumerate}

\clearpage
